\begin{problem}[Compact support is not restriction-stable]
Consider the assignment $U\mapsto C_c(U)$ of compactly supported continuous real-valued functions.  Show that ordinary restriction does not make this assignment a presheaf on arbitrary open subsets.
\end{problem}

\paragraph{Why this problem matters.}
Not every familiar space of functions can be used as presheaf sections with naive restrictions.

\paragraph{Useful ideas before starting.}
Restrict a bump supported on $[0,1]$ from $\RR$ to $(0,1)$.

\paragraph{Guided hints.}
Restrict a bump supported on $[0,1]$ from $\RR$ to $(0,1)$.

\begin{solution}
Choose $f\in C_c(\RR)$ with support $[0,1]$ and $f>0$ on $(0,1)$.  On $V=(0,1)$ the restricted function is nonzero everywhere, so its support relative to $V$ is $V$.  The space $(0,1)$ is not compact.  Thus $f|_V\notin C_c(V)$, so ordinary restriction does not define the proposed presheaf.
\end{solution}

\paragraph{What happened mathematically.}
The failure occurs before sheaf axioms: the proposed section spaces are not closed under restriction.

\paragraph{Extensions.}
Find classes of inclusions $V\subseteq U$ for which compact support is preserved.

\paragraph{Provenance.}
New canonical counterexample.
